\section{Literature Survey}
\label{sec:related}

\subsection{Prognostics and the C-MAPSS benchmark}
Prognostics estimates the \emph{remaining useful life} (RUL) of an asset from
condition-monitoring data, so that maintenance can be planned before failure
rather than after it~\cite{jardine2006review,lei2018review}. Approaches are
usually grouped into physics-based models, statistical (data-driven stochastic)
models, and machine-learning models~\cite{si2011review,lei2018review}.
The C-MAPSS dataset~\cite{saxena2008damage} was released by NASA for the 2008
PHM data challenge: turbofan engines are simulated from a healthy state to
failure under randomized initial wear, sensor noise, one or six operating
regimes and one or two fault modes. It remains the most widely used
RUL benchmark; Ramasso and Saxena~\cite{ramasso2014benchmark} surveyed the
methods applied to it and documented pitfalls in how results are reported,
including inconsistent RUL clipping and evaluation protocols.

\subsection{Statistical and model-based prognostics}
Two strands from this literature directly shape our digital twin.
\emph{Health-index (HI) construction}: Wang et~al.~\cite{wang2008similarity}
(a winning approach at PHM08) fuse sensors into a single HI by regressing
sensor values onto $1$ for early-life and $0$ for end-of-life cycles, then
match a test engine's HI curve against a library of training trajectories.
\emph{Stochastic degradation models with Bayesian updating}: Gebraeel
et~al.~\cite{gebraeel2005residual} model degradation as an exponential process
whose parameters have a population prior learned from a fleet, and update that
prior with each new observation of the individual unit to obtain a
residual-life \emph{distribution}. Particle filters~\cite{arulampalam2002tutorial}
generalize this updating to nonlinear, non-Gaussian settings; Orchard and
Vachtsevanos~\cite{orchard2009particle} established the particle-filtering
framework for online fault diagnosis and failure prognosis, and the
Liu--West kernel-shrinkage scheme~\cite{liu2001combined} is the standard way to
estimate static model parameters inside a particle filter without particle
collapse. Such models are interpretable and yield full predictive
distributions, but are seldom benchmarked against modern deep models on
C-MAPSS.

\subsection{Deep learning for RUL}
Since 2016, deep networks trained on sliding windows of raw sensor data have
dominated reported C-MAPSS accuracy. Babu et~al.~\cite{babu2016cnn} introduced
convolutional networks for RUL regression; Zheng et~al.~\cite{zheng2017lstm}
applied LSTMs~\cite{hochreiter1997lstm}; Li et~al.~\cite{li2018dcnn} showed
that deeper 1-D convolutions with a piecewise-linear RUL target reach RMSE near
12.6 on FD001. Later work added semi-supervised pre-training~\cite{listou2019semisup},
encoder--decoders~\cite{yu2019bilstmed}, hybrid CNN--LSTM
architectures~\cite{aldulaimi2019hdnn} and attention~\cite{ragab2020ats2s}.
Table~\ref{tab:literature} compiles representative results. Two gaps are
consistent across this line of work: (i) models output a point estimate,
with no calibrated measure of confidence, and (ii) accuracy is reported only at
the last observed cycle of each truncated test engine, which says nothing about
whether the resulting \emph{maintenance decisions} are safe or economical.
Gradient-boosted trees~\cite{friedman2001gbm,ke2017lightgbm} on hand-crafted
window features remain a strong, far cheaper baseline.

\subsection{Uncertainty quantification}
Monte-Carlo dropout~\cite{gal2016dropout} and deep
ensembles~\cite{lakshminarayanan2017ensembles} are the most common ways to
attach uncertainty to neural RUL predictors. Their calibration is not
guaranteed: Basora et~al.~\cite{basora2023uqbench} benchmark these methods for
prognostics and find that calibration quality varies substantially across
methods and engines. Conformal
prediction~\cite{vovk2005alrw,angelopoulos2021gentle} instead wraps any point
predictor and provides intervals with a finite-sample, distribution-free
coverage guarantee under exchangeability; Mondrian (group-conditional) variants
calibrate separately within bins, which matters for RUL where error grows with
distance from failure. To our knowledge conformal intervals are still rarely
used to drive maintenance actions on C-MAPSS.

\subsection{Digital twins}
The digital-twin concept was articulated for aerospace by Glaessgen and
Stargel~\cite{glaessgen2012digital} as a model that ``mirrors the life of its
flying twin'' by combining physics, fleet history and the vehicle's own sensor
data, and was formalized by Grieves and Vickers~\cite{grieves2017digital} as a
virtual counterpart linked to one specific physical instance and kept
synchronized with it. Tao et~al.~\cite{tao2019digital} review industrial uses,
with prognostics and health management among the most prominent. Two
properties recur in these definitions and are often missing from works that
label a neural RUL regressor a ``digital twin'': the model is
\emph{instance-specific} (it holds state about one asset) and it is
\emph{synchronized} with that asset as data arrive.

\subsection{Edge intelligence and cascaded inference}
Edge computing places computation near data sources to reduce latency,
bandwidth and dependence on connectivity~\cite{shi2016edge}. Cascaded or
early-exit inference~\cite{teerapittayanon2016branchynet} runs a cheap model on
every input and defers to an expensive one only when the cheap model is not
confident enough. Our framework applies the same idea across a
device/twin boundary: a stateless edge model runs every cycle, and the
stateful twin is synchronized only when a trigger fires.

\subsection{From predictions to maintenance decisions}
Classical reliability theory gives the baseline any predictive policy must
beat: \emph{age replacement}, which replaces a unit at a fixed age chosen to
minimize the long-run cost rate of the renewal process~\cite{barlow1960optimum}.
We evaluate every policy with the same renewal-reward cost rate, so predictive
maintenance is compared against this optimal time-based policy and not only
against run-to-failure.

\subsection{Positioning of this work}
Existing C-MAPSS studies optimize last-cycle RMSE of point predictors;
existing digital-twin work is mostly conceptual or not benchmarked on it.
This project combines (i) an instance-specific, Bayesian-updated degradation
twin in the tradition of~\cite{gebraeel2005residual,orchard2009particle},
(ii) a conformally calibrated edge model, and (iii) a trigger-based
synchronization scheme, and evaluates the combination on what matters
operationally: failures prevented, useful life used and cost rate, over
engines replayed all the way to failure.
